\chapter{Requirements Analysis}

\section{Introduction}

This chapter presents the requirements analysis of the proposed system. The goal of this analysis is to identify the main functional and non-functional requirements needed to develop an automation platform that allows users to create, manage, and execute workflows through a visual interface. The system provides integration with external services through APIs and supports automated execution of different tasks using predefined nodes.

\section{System Overview}

The system is an automation platform that enables users to build workflows by connecting different nodes together. Each node represents a specific operation, such as processing data, communicating with external APIs, or performing a computational task. The created workflows are stored and executed by the backend execution engine.

The platform also provides AI-assisted workflow generation, allowing users to describe their desired automation process using natural language. The system analyzes the request and generates a suitable workflow structure using specialized AI components.

\section{Functional Requirements}

Functional requirements describe the main features and operations provided by the system.

\subsection{User Interface}

The system shall provide a graphical interface that allows users to:

\begin{itemize}
    \item Create and edit workflows using a visual node-based editor.
    \item Add, remove, and connect workflow nodes.
    \item Configure node parameters.
    \item Save and load existing workflows.
\end{itemize}

\subsection{Workflow Management}

The system shall support workflow management operations, including:

\begin{itemize}
    \item Creating new workflows.
    \item Storing workflow definitions.
    \item Retrieving previously created workflows.
    \item Updating workflow configurations.
    \item Executing selected workflows.
\end{itemize}

\subsection{Node Management}

The system shall provide a collection of predefined nodes that represent different operations.

Each node shall contain:

\begin{itemize}
    \item A unique identifier.
    \item A description of its functionality.
    \item Required input parameters.
    \item Execution logic.
    \item Defined input and output formats.
\end{itemize}

The system shall allow the execution engine to discover and execute nodes dynamically through a node registry.

\subsection{API Integration}

The system shall support communication with external services through APIs.

The API layer shall provide:

\begin{itemize}
    \item Endpoints for workflow creation and management.
    \item Endpoints for workflow execution.
    \item Communication between the frontend application and backend services.
    \item Integration with external services used by workflow nodes.
\end{itemize}

\subsection{Workflow Execution}

The system shall provide an execution engine responsible for running workflows.

The execution engine shall:

\begin{itemize}
    \item Receive workflow execution requests.
    \item Process workflow nodes in the correct order.
    \item Manage execution tasks asynchronously.
    \item Handle communication between connected nodes.
    \item Report execution results and errors.
\end{itemize}

\subsection{AI-Assisted Workflow Generation}

The system shall provide an AI assistant that helps users generate workflows from natural language descriptions.

The AI workflow generation module shall:

\begin{itemize}
    \item Analyze user requests.
    \item Identify required workflow operations.
    \item Search available nodes.
    \item Construct a workflow representation.
    \item Validate the generated workflow before execution.
\end{itemize}

\section{Non-Functional Requirements}

Non-functional requirements describe the quality attributes of the system.

\subsection{Performance}

The system should:

\begin{itemize}
    \item Execute workflows efficiently.
    \item Handle multiple workflow execution requests.
    \item Minimize delays between workflow steps.
\end{itemize}

\subsection{Scalability}

The system should be designed to support:

\begin{itemize}
    \item Adding new node types without modifying existing components.
    \item Increasing the number of available workflows and users.
    \item Extending integrations with additional external APIs.
\end{itemize}

\subsection{Reliability}

The system should:

\begin{itemize}
    \item Handle execution failures gracefully.
    \item Provide meaningful error messages.
    \item Maintain consistent workflow data.
\end{itemize}

\subsection{Maintainability}

The system should have a modular architecture that allows:

\begin{itemize}
    \item Independent development of system components.
    \item Easy addition of new features.
    \item Clear separation between workflow management, execution, and AI components.
\end{itemize}

\section{Summary}

This chapter identified the main requirements of the proposed automation platform. The system focuses on providing workflow creation, API-based communication, dynamic node execution, and AI-assisted workflow generation while maintaining a modular and extensible architecture.